\chapter{系统实现}
\label{chap:implementation}

本章说明系统的开发环境、项目结构和关键模块的实现。重点不在逐行解释代码，而在说明每个模块如何落实第~\ref{chap:design} 章的设计：认证与授权如何拦截请求，文档如何被解析和分块，双路召回如何融合，提示词如何约束生成，引用如何校验，统计如何聚合。代码片段只保留与论文设计对应的核心部分，完整的配置、异常处理和前端页面不在正文中展开。

\section{开发环境与项目结构}
\label{sec:environment}

系统采用前后端分离结构，开发环境如表~\ref{tab:env} 所示。后端使用 Python 3.11，前端使用 Node.js 20，数据库为 MySQL 8.0，向量库为 Chroma，嵌入模型选用 BGE 中文小模型，重排模型选用对应的交叉编码器，大模型通过兼容 OpenAI 协议的 HTTP 接口调用。开发工具包括 VS Code、Navicat、Postman 和 Git。

\begin{table}[H]
  \centering
  \caption{开发环境}
  \label{tab:env}
  \zihao{5}
  \renewcommand{\arraystretch}{1.25}
  \begin{tabular}{p{3.0cm}p{9.6cm}}
    \toprule
    类别 & 配置 \\
    \midrule
    操作系统 & Windows 11，开发与演示环境 \\
    后端 & Python 3.11，FastAPI，Uvicorn，SQLAlchemy \\
    前端 & Node.js 20，Vue 3，Vite，Element Plus，Axios \\
    数据库 & MySQL 8.0，InnoDB，utf8mb4 \\
    向量库 & Chroma，本地持久化模式 \\
    嵌入与重排 & BGE 中文小模型，交叉编码器重排模型 \\
    大模型 & 兼容 OpenAI 协议的可配置 API \\
    文档解析 & PyMuPDF、python-docx、chardet \\
    版本管理 & Git，GitHub 远程仓库 \\
    \bottomrule
  \end{tabular}
\end{table}

项目分为后端、前端和部署三个部分。后端按领域划分目录：\texttt{api} 存放路由，\texttt{core} 存放配置、安全和依赖，\texttt{models} 存放 SQLAlchemy 模型，\texttt{schemas} 存放请求与响应结构，\texttt{services} 存放认证、文档、检索、生成和统计服务，\texttt{tasks} 存放异步任务，\texttt{utils} 存放分块、文本清洗和引用解析工具。前端按页面和组件划分：\texttt{views} 存放登录、问答、历史、资料、统计和配置页面，\texttt{components} 存放消息气泡、引用卡片、上传组件和图表组件，\texttt{api} 封装接口请求，\texttt{stores} 管理用户和会话状态。部署部分包括 Nginx 配置、环境变量示例、数据库初始化脚本和启动说明。

配置通过环境变量注入，包括数据库连接、JWT 密钥、大模型地址与密钥、嵌入模型路径、向量库路径、检索参数和日志级别。仓库只提交 \texttt{.env.example}，不提交真实密钥；本地开发使用 \texttt{.env}，部署时由运行环境注入。这样做可以避免密钥进入版本历史，也便于在不同环境之间切换。

\section{用户与权限管理}
\label{sec:auth-impl}

登录接口接收用户名和密码，先根据用户名查询用户记录，再使用 bcrypt 校验密码。校验通过后，服务读取用户的角色和权限，生成 JWT。令牌中只放用户编号、角色编号、签发时间和过期时间，不放密码、手机号等敏感信息。前端把令牌保存在内存或 \texttt{sessionStorage} 中，每次请求由 Axios 拦截器添加到 \texttt{Authorization} 头；令牌过期时，前端跳转登录页，避免在页面中反复弹出错误。

后端的认证与授权通过 FastAPI 依赖注入实现。认证依赖解析令牌、校验签名和有效期、查询用户状态；授权依赖检查当前用户是否具备路由声明的权限。核心代码如下所示。

\begingroup\linespread{1.0}\selectfont
\begin{lstlisting}[language=Python]
def get_current_user(
    token: str = Depends(oauth2_scheme),
    db: Session = Depends(get_db),
) -> User:
    payload = decode_token(token)
    user = db.get(User, int(payload["sub"]))
    if user is None or user.status != "active":
        raise HTTPException(status_code=401, detail="invalid token")
    return user


def require_permission(permission: str):
    def checker(user: User = Depends(get_current_user)) -> User:
        if permission not in user.permission_codes:
            raise HTTPException(status_code=403, detail="forbidden")
        return user
    return checker
\end{lstlisting}
\endgroup

路由上通过 \texttt{Depends(require\_permission("document:write"))} 声明权限，业务代码不需要重复判断角色。对于“只能查看自己的会话”这类数据级权限，服务层在查询时附加 \texttt{user\_id} 条件，而不是先查出全部数据再在前端过滤。管理员接口还会记录操作日志，包括操作人、操作对象、操作类型和结果，便于审计。

前端根据登录接口返回的权限集合生成菜单。没有权限的菜单不渲染，按钮置灰或隐藏；但前端隐藏只是体验优化，真正的拦截在后端。令牌刷新采用双令牌机制：访问令牌过期后，前端用刷新令牌换取新的访问令牌；刷新令牌也过期时，跳转登录页。注销时，前端清除本地令牌，后端把当前令牌加入短期失效列表，避免令牌在有效期内被继续使用。

\section{文档管理与知识库构建}
\label{sec:doc-impl}

文档上传接口接收 multipart 表单，字段包括文件、标题、分类、权限标签和备注。服务先计算文件 SHA-256，查询是否已有相同哈希的文档；若存在，返回重复提示，由用户选择覆盖或跳过。随后校验扩展名、MIME 类型、大小和页数，保存原文件到按日期分层的目录，并在 \texttt{document} 表插入记录，状态为“待处理”。接口立即返回文档编号，解析任务交给后台执行。

解析服务根据文件类型选择解析器。PDF 使用 PyMuPDF 逐页提取文本和页码，DOCX 使用 python-docx 读取段落、表格和样式，TXT 使用 chardet 探测编码后读取。解析结果统一为“块”列表，每个块包含文本、页码、样式层级和类型（标题、正文、表格、列表）。清洗时去除页眉页脚、连续空行和重复字符，保留条款编号、日期和金额等关键信息。

分块服务先按标题层级把块序列切成章节，再在章节内按段落和句子累积，直到达到目标长度。若当前段落加入后超过最大长度，则在句子边界处截断；相邻片段保留一定重叠，避免答案被切断。核心逻辑如下所示。

\begingroup\linespread{1.0}\selectfont
\begin{lstlisting}[language=Python]
def structure_aware_chunks(blocks, max_len=800, overlap=100):
    chunks, buf, heading = [], [], ""
    for block in blocks:
        if block.kind == "heading":
            if buf: chunks.append(make_chunk(buf, heading)); buf = []
            heading = block.text.strip(); continue
        if len("".join(buf)) + len(block.text) <= max_len:
            buf.append(block.text)
        else:
            chunks.append(make_chunk(buf, heading))
            buf = ["".join(buf)[-overlap:], block.text]
    if buf: chunks.append(make_chunk(buf, heading))
    return chunks
\end{lstlisting}
\endgroup

每个片段写入 \texttt{chunk} 表，保存文档编号、原文、页码、标题路径和向量编号；同时调用嵌入模型生成向量，写入 Chroma。写入采用“先删旧、再写新”的策略：重建某篇文档时，先按文档编号删除旧片段和旧向量，再批量写入新片段；如果中途失败，任务状态记为失败，旧索引保持不变，避免知识库出现不完整的中间状态。批量写入使用固定大小的批次，降低内存占用；嵌入模型不可用时，任务进入重试队列，最多重试三次，仍失败则通知管理员。

知识库构建页面展示文档列表、处理状态、片段数量、更新时间和操作按钮。管理员可以按分类、状态和关键词筛选，查看失败原因，单篇重试或批量重建。前端轮询任务状态，处理中显示进度条，成功后提示片段数量，失败时展开错误详情。对于扫描版 PDF，解析结果为空时提示“该文件可能需要 OCR”，并给出转为文本后重新上传的建议。

\section{混合检索与重排序实现}
\label{sec:retrieval-impl}

检索服务对外提供统一的 \texttt{search(query, user, top\_k)} 接口，内部依次执行 BM25 召回、向量召回、RRF 融合和交叉编码器重排。BM25 使用中文分词后的词项建立索引，索引在文档入库时增量更新；向量召回调用嵌入模型编码查询，再在 Chroma 中按余弦相似度取前 $k_1$ 条。两路召回都带上权限过滤条件，确保候选片段对当前用户可见。

RRF 融合只依赖排名，不依赖原始分数，因此不需要对 BM25 分数和余弦相似度做归一化。实现时把两路召回结果转换为“片段编号列表”，再按式~\eqref{eq:rrf} 累加分数。核心代码如下所示。

\begingroup\linespread{1.0}\selectfont
\begin{lstlisting}[language=Python]
def reciprocal_rank_fusion(rankings, k=60):
    scores = defaultdict(float)
    for ranking in rankings:
        for rank, chunk_id in enumerate(ranking, start=1):
            scores[chunk_id] += 1.0 / (k + rank)
    return sorted(scores.items(), key=lambda x: x[1], reverse=True)
\end{lstlisting}
\endgroup

融合后取前 $k_2$ 条候选送入重排模型。重排模型把查询和片段拼接后编码，输出相关性分数；服务按分数降序排序，丢弃低于阈值 $\tau$ 的候选。若候选为空或全部低于阈值，返回空结果，由问答服务执行拒答。重排阶段设置超时，超时则退化为 RRF 排序，保证接口在模型服务不稳定时仍能返回结果。

为了观察检索效果，服务记录每次请求的查询、召回通道、候选数量、融合耗时、重排耗时和最终分数。日志按请求编号关联，便于分析“某类问题总是召回不到”或“重排把正确片段排到后面”的原因。管理端提供检索调试页面，管理员输入问题后可以查看两路召回结果、融合排名和重排分数，辅助调整参数和补充资料。

\section{受约束生成、引用溯源与多轮会话}
\label{sec:generation-impl}

生成服务把系统提示、证据、问题和输出约束组装成消息列表，调用大模型接口。系统提示明确角色和边界；证据按编号排列，每条包含标题、章节路径和原文；用户问题保留原始表述；输出约束要求先给结论、再给依据、关键结论标注引用编号、不得编造资料中没有的信息、证据不足时说明无法回答。提示词模板的核心部分如下所示。

\begingroup\linespread{1.0}\selectfont
\begin{lstlisting}[language=Python]
SYSTEM_PROMPT = """You are a campus question-answering assistant.
Answer ONLY with the evidence below. Do not use outside knowledge.
Cite the evidence number in square brackets after each key claim.
If the evidence is insufficient, say you cannot answer and explain why.
"""

USER_PROMPT = """[Evidence]
{evidence}

[Question]
{question}

[Output]
1. A concise answer with citations like [1].
2. A short list of the evidence numbers you used.
"""
\end{lstlisting}
\endgroup

生成结果先经过引用解析。服务用正则表达式提取答案中的 \texttt{[n]} 编号，检查编号是否在候选证据集合内；再检查每个编号对应的片段是否存在、当前用户是否有权限查看；最后把答案中的关键句与引用片段做关键词覆盖检查。若引用编号越界或片段无效，服务隐藏无效引用，并把答案降级为摘录式回答；若答案没有任何引用，服务在末尾追加“以上内容未找到明确依据”，并标记为需人工检查。

多轮会话在生成前执行问题改写。服务读取最近 $N$ 轮消息，把当前问题、上一轮问题和上一轮答案送入改写模型，要求输出一个不依赖上下文的自包含问题。改写结果用于检索，原始问题用于生成和日志。若改写结果为空或与原问题差异过大，服务保留原始问题作为兜底查询，两路查询分别召回后合并，避免改写错误导致完全跑题。会话消息保存问题时同时保存改写后的问题、引用片段和耗时，便于后续分析。

拒答策略根据证据情况分三级。第一级，检索为空或最高重排分数低于阈值，直接返回“暂时找不到相关依据”，并建议用户换一种问法或选择分类。第二级，证据存在但引用校验失败，返回检索到的原文片段，并说明“以下为检索到的资料，未生成总结”。第三级，模型超时或返回格式错误，重试一次；仍失败则返回原文片段，并记录失败日志。三级策略的共同目标是：宁可少答，也不返回没有依据的内容。

\section{管理统计与系统配置}
\label{sec:stats-impl}

统计服务按时间范围、分类和用户角色聚合问答数据，提供高频问题、未解决问题、知识库使用量、满意度和平均响应时延等指标。高频问题按规范化后的问题文本分组，去除标点和停用词，避免“补办学生证需要什么材料”和“补办学生证要哪些材料”被算成两个问题；未解决问题包括拒答、引用校验失败和用户点踩的消息；知识库使用量统计文档数量、片段数量、检索次数和命中片段分布；满意度按好评率和平均分展示；响应时延按检索、重排、生成三段拆分，便于定位瓶颈。

统计接口使用分页和聚合查询，避免一次性加载全部明细。高频问题按出现次数降序取前 20 条，未解决问题按时间倒序取前 50 条，趋势图按天聚合。管理端使用表格和图表展示，支持导出 CSV。对于数据量较大的场景，可以增加按天预聚合的统计表，定时任务把明细汇总为日指标，查询时直接读取汇总表。

系统配置页面提供检索条数、RRF 常数、重排阈值、模型名称、超时时间、拒答提示语等配置项。配置保存前做类型和范围校验，例如 $k_1$ 必须在 5 到 100 之间，$\tau$ 必须在 0 到 1 之间，超时时间必须在 3 到 30 秒之间；保存后写入 \texttt{config} 表并刷新缓存，下一次请求生效。配置修改记录操作人和修改前后值，便于回滚。密钥类配置不在页面展示，只允许通过环境变量注入，避免管理员误操作或页面泄露。

